\scanpage{7}

\uline{\illegible{1}}\unsure{:|,} \unsure{soit} $\unsure{\mathcal{A}|\mathscr{A}|A}$ un anneau sur $X$, soient $M,N$ respectivement les faisceaux $K^{0}$ et $K^{1}$ associés au champ additif des $\mathcal{A}$-Modules projectifs de \unsure{v.|t.|r.} f. \inserted{(p ex.),} Alors la \unsure{construction} de Mme Sinh nous fourni un champ de Picard \unsure{e|é}pinglé par $M,N$, d'où une section \deleted{\unsure{du|de}} canonique du terme médian $\unsure{P|F}(M,N$ de $(*)$.

\uline{NB} Tout ce qui précède a les \unsure{fonctorialités} évidentes en $M,N,X$\ldots

\uline{Question} : Le triangle exact $(\unsure{T|J})$ \unsure{et} la suite exacte $(\unsure{*})$ sont ils \unsure{connu|connus} par les comp\unsure{è|é}tences (Quillen, Breen Illusie\ldots) ? \unsure{Connaissent-ils|Donnaissent-ils} des variantes ``supérieures'' ? (Un principe ``géométrique'' pour en obtenir pourrait être via des $n$-champs de Picard non néc. stricts\ldots)

% Two curved impressions run along the lower left edge of the copied sheet.
% archive sheet number: 64
